%%%PAGE 267 rot=0 legibility=good summary=等效场中圆周运动求最大拉力；摩擦角法；带电粒子在磁场中的动量定理；x-t图像折返问题；斜面与水平面上的追及问题
%%%BLOCK id=267-01 ch=04 sec=竖直面圆周运动·等效重力场 kind=problem alt=09 cont=none y=0.00-0.25
\begin{problem}
%FIG p267_a 0.07 0.02 0.36 0.18
\notefig[0.4]{figures/p267_a.png}
（图注：等效最低点）\quad $E=\dfrac{mg}{q}$\qquad 求 $F_{\max}$
\begin{solution}
\[
\left\{\begin{aligned}
&\sqrt2\,mg\sqrt2\,R=\dfrac12mv^2\\
&F-\sqrt2\,mg=\dfrac{mv^2}{R}\\
&\sqrt2\,mgR\left(1-\dfrac{\sqrt2}{2}\right)=\dfrac12m\left(v'^2-v_x^2\right)
\end{aligned}\right.
\]
% CHECK: 第三式括号内 v'^2 的撇号/指数手写不清
\end{solution}
\end{problem}
%%%END
%%%BLOCK id=267-02 ch=90 sec=解题注意 kind=note alt=none cont=none y=0.265-0.31
单位\unclear{$\ne$}物理量\qquad 求比值要求绝对值
%%%END
%%%BLOCK id=267-03 ch=02 sec=共点力平衡·摩擦角法 kind=method alt=none cont=none y=0.31-0.48
\textbf{使用摩擦角}
%FIG p267_b 0.02 0.305 0.33 0.485
\notefig[0.45]{figures/p267_b.png}
$f\downarrow$\quad $F$ 先 $\downarrow$ 后 $\uparrow$.
%%%END
%%%BLOCK id=267-04 ch=11 sec=洛伦兹力·动量定理 kind=problem alt=06 cont=none y=0.48-0.60
\begin{problem}
%FIG p267_c 0.03 0.478 0.255 0.59
\notefig[0.3]{figures/p267_c.png}
\begin{solution}
\[ -BqV_xt-mgt=0-mV_0 \]
$V_{\text{末}}$ 用动能定理求
\[ t=\dfrac{mV_0-qB\dfrac{L}{2}}{mg} \]
\end{solution}
\end{problem}
%%%END
%%%BLOCK id=267-05 ch=01 sec=图像·x-t图像 kind=method alt=03 cont=none y=0.60-0.77
%FIG p267_d 0.04 0.605 0.30 0.73
\notefig[0.35]{figures/p267_d.png}
（图中标注：$x_0$、$V_1$、$V_2$、$T$、$2T$，前段标“一次函数”，后段标“二次函数”）

$0\sim T$\quad $F_{\text{合}}$ 为 $0$

匀加速\quad \unclear{a}变了方向 故 $F$ 不为恒力
% CHECK: “匀加速”后的单字手写不清（像 a/t/↑），“方向”前的主语待核
\[ -V_1T=\dfrac{-V_1+V_2}{2}T\qquad\text{位移大小相等方向相反} \]
% CHECK: 按“位移大小相等方向相反”及结论 V_2=3V_1，左边应为 V_1T（原稿左边多一个负号，或右边应取负）
\[ \therefore\ V_2=3V_1 \]
%%%END
%%%BLOCK id=267-06 ch=01 sec=追及相遇·斜面与水平面 kind=problem alt=03 cont=none y=0.77-1.00
\begin{problem}
%FIG p267_e 0.04 0.78 0.405 0.885
\notefig[0.45]{figures/p267_e.png}
在 C 处甲刚好追上乙
\begin{solution}
$t_1=t_2\ \to\ 1:3:\cdots$\qquad $S_{\text{乙}}=\dfrac14BC$
% CHECK: S 的下标手写像 D/2，按上下文取“乙”
\[
\begin{matrix}\text{甲：}\\ \text{乙：}\end{matrix}
\left\{\begin{aligned}
X_2&=Vt_2\\
X_2&=\dfrac{V}{2}(t_1+t_2)\qquad t_1=t_2
\end{aligned}\right.
\]
\[ \text{甲：}\ X_1=\dfrac{V}{2}t_1=\dfrac12Vt_2=\dfrac12X_2\qquad \therefore\ AB=\dfrac12BC \]
\[ X_1=\dfrac12at_1^2\qquad X_2=\dfrac12\times\dfrac14g(t_1+t_2)^2=\dfrac12gt_1^2\qquad \therefore\ a=\dfrac12g\qquad(\theta=30^\circ) \]
\end{solution}
\end{problem}
%%%END
%%%PAGEEND 267
